% ====================================================================
% Файл: tasks/01_mechanics/kim06.tex
% Тема: Задание №6 (Изменение величин и соответствие графиков в кинематике/динамике)
% Контекст: Свободное падение и движение тела, брошенного под углом к горизонту
% ====================================================================

% === ЗАДАЧА 11 ===
% Тег: #МЕХАНИКА #КИМ_06 #ВАРИАНТ_11
\begin{taskbasic}[code={\label{task:dem26_v11_n6}}]
\textbf{Задача 11.} В момент $t=0$ камень бросают с начальной скоростью $\vec{v}_0$ под углом $\alpha$ к горизонту с некоторой высоты $h$. Графики А и Б представляют собой зависимости физических величин, характеризующих движение камня в процессе полёта, от времени $t$. 

Установите соответствие между графиками и физическими величинами, зависимость которых от времени эти графики могут представлять. (Сопротивлением воздуха пренебречь. Потенциальная энергия камня отсчитывается от уровня $y=0$.)

\nopagebreak\smallskip
\begin{center}
\begin{minipage}[b]{0.45\linewidth}
\centering
\textbf{ГРАФИКИ}\\
\smallskip
\begin{tikzpicture}[scale=0.7, every node/.style={font=\footnotesize}]
    \node[above] at (1.5,2.5) {А)};
    \draw[->] (0,0) -- (3.5,0) node[right] {$t$};
    \draw[->] (0,0) -- (0,2.5) node[above] {};
    \draw[ultra thick, brandPrimary, domain=0:3, samples=100] plot (\x, {1.2 + 0.8*\x - 0.4*\x*\x});
    \node[below left] at (0,0) {0};
\end{tikzpicture}

\vspace{0.5cm}
\begin{tikzpicture}[scale=0.7, every node/.style={font=\footnotesize}]
    \node[above] at (1.5,1.5) {Б)};
    \draw[->] (0,0) -- (3.5,0) node[right] {$t$};
    \draw[->] (0,0) -- (0,2.0) node[above] {};
    \draw[ultra thick, brandPrimary] (0,1.2) -- (3,1.2);
    \node[below left] at (0,0) {0};
\end{tikzpicture}
\end{minipage}
\begin{minipage}[b]{0.5\linewidth}
\textbf{ФИЗИЧЕСКИЕ ВЕЛИЧИНЫ}\\
\smallskip
1) потенциальная энергия камня \\
2) проекция импульса камня на ось $x$ \\
3) проекция ускорения камня на ось $y$ \\
4) кинетическая энергия камня \\
\end{minipage}
\end{center}

К каждой позиции первого столбца подберите соответствующую позицию из второго столбца и запишите в таблицу выбранные цифры под соответствующими буквами.

\kim{6}
\end{taskbasic}
\answer{12}

% === ЗАДАЧА 12 ===
% Тег: #МЕХАНИКА #КИМ_06 #ВАРИАНТ_12
\begin{taskbasic}[code={\label{task:dem26_v12_n6}}]
\textbf{Задача 12.} В момент $t=0$ камень бросают с начальной скоростью $\vec{v}_0$ под углом $\alpha$ к горизонту с некоторой высоты $h$. Графики А и Б представляют собой зависимости физических величин, характеризующих движение камня в процессе полёта, от времени $t$.

Установите соответствие между графиками и физическими величинами, зависимость которых от времени эти графики могут представлять. (Сопротивлением воздуха пренебречь. Потенциальная энергия камня отсчитывается от уровня $y=0$.)

\nopagebreak\smallskip
\begin{center}
\begin{minipage}[b]{0.45\linewidth}
\centering
\textbf{ГРАФИКИ}\\
\smallskip
\begin{tikzpicture}[scale=0.7, every node/.style={font=\footnotesize}]
    \node[above] at (1.5,2.5) {А)};
    \draw[->] (0,0) -- (3.5,0) node[right] {$t$};
    \draw[->] (0,0) -- (0,2.5) node[above] {};
    \draw[ultra thick, brandPrimary, domain=0:3, samples=100] plot (\x, {1.5 - 0.8*\x + 0.4*\x*\x});
    \node[below left] at (0,0) {0};
\end{tikzpicture}

\vspace{0.5cm}
\begin{tikzpicture}[scale=0.7, every node/.style={font=\footnotesize}]
    \node[above] at (1.5,1.5) {Б)};
    \draw[->] (0,0) -- (3.5,0) node[right] {$t$};
    \draw[->] (0,-1.5) -- (0,1.5) node[above] {};
    \draw[ultra thick, brandPrimary] (0,-1.0) -- (3,-1.0);
    \node[below left] at (0,0) {0};
\end{tikzpicture}
\end{minipage}
\begin{minipage}[b]{0.5\linewidth}
\textbf{ФИЗИЧЕСКИЕ ВЕЛИЧИНЫ}\\
\smallskip
1) потенциальная энергия камня \\
2) проекция импульса камня на ось $y$ \\
3) проекция ускорения камня на ось $y$ \\
4) кинетическая энергия камня \\
\end{minipage}
\end{center}

К каждой позиции первого столбца подберите соответствующую позицию из второго столбца и запишите в таблицу выбранные цифры под соответствующими буквами.

\kim{6}
\end{taskbasic}
\answer{43}

% === ЗАДАЧА 23 ===
% Тег: #МЕХАНИКА #КИМ_06 #ВАРИАНТ_23
\begin{taskbasic}[code={\label{task:dem26_v23_n6}}]
\textbf{Задача 23.} Шарик, брошенный горизонтально с высоты $H$ с начальной скоростью $\vec{v}_0$, за время $t$ пролетел в горизонтальном направлении расстояние $L$.

Что произойдёт со временем полёта и дальностью полёта шарика, если на этой же установке уменьшить начальную скорость шарика в $2$ раза? Сопротивлением воздуха пренебречь.

Для каждой величины определите соответствующий характер её изменения:
1) увеличится;
2) уменьшится;
3) не изменится.

Запишите в таблицу выбранные цифры для каждой физической величины. Цифры в ответе могут повторяться.
\begin{center}
\begin{tabular}{|c|c|}
\hline
Время полёта & Дальность полёта \\ \hline
 & \\ \hline
\end{tabular}
\end{center}

\kim{6}
\end{taskbasic}
\answer{32}

% === ЗАДАЧА 24 ===
% Тег: #МЕХАНИКА #КИМ_06 #ВАРИАНТ_24
\begin{taskbasic}[code={\label{task:dem26_v24_n6}}]
\textbf{Задача 24.} Шарик, брошенный горизонтально с высоты $H$ с начальной скоростью $\vec{v}_0$, за время $t$ пролетел в горизонтальном направлении расстояние $L$.

Что произойдёт со временем полёта и дальностью полёта шарика, если на этой же установке уменьшить высоту $H$ в $2$ раза, а начальную скорость оставить прежней? Сопротивлением воздуха пренебречь.

Для каждой величины определите соответствующий характер её изменения:
1) увеличится;
2) уменьшится;
3) не изменится.

Запишите в таблицу выбранные цифры для каждой физической величины. Цифры в ответе могут повторяться.
\begin{center}
\begin{tabular}{|c|c|}
\hline
Время полёта & Дальность полёта \\ \hline
 & \\ \hline
\end{tabular}
\end{center}

\kim{6}
\end{taskbasic}
\answer{22}